By construction measurements in an effective field theory take place at energies below the cutoff scale of the EFT.
As such, to link measurements to parameters of possible physics in the UV, it is crucial to understand how the coefficients of the EFT operators evolve from one scale to another.
This behavior is described by the renormalization group evolution equations
\begin{equation}
\gamma_{ij} C_j = 16 \pi^2 \mu \frac{d C_i}{d \mu} \equiv \dot C_i ,
\end{equation}
where $\gamma$ is the anomalous dimension matrix and $\mu$ is the renormalization scale.

The renormalization group evolution of the SMEFT operators at one-loop is known at dimension-5~\cite{Babu:1993qv}, dimension-6~\cite{Jenkins:2013zja, Jenkins:2013wua, Alonso:2013hga, Alonso:2014zka}, and dimension-7~\cite{Liao:2016hru, Liao:2019tep}.
The loop diagrams in the preceding computations were built from one dimension-$d$ SMEFT contact amplitude and one tree level SM amplitude.
For the dimension-6 RGE there are also one-loop amplitudes resulting from two insertions of dimension-5 operators, and Ref.~\cite{Davidson:2018zuo} computed these contributions to the RGE equations.
Typically, the two dimension-5 operator contribution to the dimension-6 RGE will be suppressed with respect to the one dimension-6 contribution because all odd mass dimension operators in the SMEFT violate at least one of baryon and lepton number~\cite{Kobach:2016ami, Helset:2019eyc}.

It is well beyond the scope of this work to compute the RGE equations for the dimension-8 operators.
Instead we content ourselves to make a few comments about the structure of the associated anomalous dimension matrix. 
The renormalization of a dimension-8 operator can happen due to one insertion of a dimension-8 operator in a loop amplitude, two insertions of dimension-6 operators, or one insertion of dimension-5 operator and one insertion of a dimension-7 operator.
The lattermost type of amplitude will typically be suppressed as was the case with the two dimension-5 operator contribution to the dimension-6 RGE.
However, the two dimension-6 operator contribution to the dimension-8 RGE will generally be comparable in magnitude to the one dimension-8 operator contribution.
Dimension-8 is the lowest mass dimension where there are two co-leading contributions to the RGE.

Many of the entries in the anomalous dimension matrix vanish~\cite{Cheung:2015aba} including beyond one-loop~\cite{Bern:2020ikv}.
One way to understand these zeroes in the anomalous dimension matrix is through non-renormalization theorem of Ref.~\cite{Cheung:2015aba}, which applies to loop amplitudes with one dimension-$d$ SMEFT insertion and one SM tree level amplitude.
To start define holomorphic and anti-holomorphic weights as $w = n - h$ and $\overline w = n + h$, respectively, where $n$ is the number of particles created by an operator and $h$ is the sum of the helicities of the particles created.
Then the theorem can be stated as operators belonging to subclass $j$ can renormalize operators of another subclass $i$ if $w_i \geq w_j$ and $\overline w_i \geq \overline w_j$.
The theorem can only be violated when the SM tree amplitude contains two Yukawa interactions.
The results of the non-renormalization theorem applied to dimension-8 operators are visualized in Table~\ref{tab:rge}.
Visually the theorem transitions down or to the left in the weight lattice.
Furthermore, if the SMEFT is UV-completed by a weak-coupled and renormalizable theory, then only the operator classes in red can be generated via tree level matching, whereas the black and blue operator classes are only generated at loop level in this scenario.
However, the operator classes in blue can be renormalized by the red tree level operator classes even though the blue operator classes are not generated by these types of UV theories at tree level~\cite{Craig:2019wmo}.

A contact amplitude resulting from an insertion of dimension-$d$ SMEFT operator (with $d$ even) obeys the relation
\begin{equation}
\label{eq:smeft_n}
2 d \geq w_k + \overline w_k \geq d .
\end{equation}
Furthermore, restricting to particles of spin-1 or less we also have
\begin{equation}
\label{eq:smeft_h}
d \geq |w_k - \overline w_k| .
\end{equation}
The relations~\eqref{eq:smeft_n} and~\eqref{eq:smeft_h} do not, in general, respect the weight bounds of SM tree level amplitudes, $w_k \geq 4$ and $\overline w_k \geq 4$ (which apply to non-exceptional amplitudes).
Therefore the theorem needs to be modified to determine vanishing entries among the two dimension-6 contributions to the dimension-8 anomalous dimension matrix as the theorem, as currently formulated, relies on the SM tree level weight bounds.
It would be interesting to examine in detail the structure of the anomalous dimension matrix resulting from two insertion of dimension-6 operators.

Inverting the logic of the non-renormalization theorem, operators of the same type will generically mix under RGE evolution.
An interesting restriction on how general the mixing can be comes from Lagrangian terms with multiple flavor structures.
Examples of such terms include $Q_{l^4H^2}$, $Q_{q^4H^2}^{(1,3)}$, $Q_{u^4H^2}$ and $Q_{d^4H^2}$.
The different flavor structures in these terms can only be mixed by Yukawa contributions to the RGE, \textit{i.e.} gauge and $\lambda$ contributions do not mix different flavor structures of a given term.
See Ref.~\cite{Alonso:2014zka} for further discussion on this point.
On the other hand, at the dimension-6, there are plenty of example of mixing amongst different terms of the same type.
For example, the gauge contribution to the dimension-6 RGE mixes the bosonic operators $C_{H\Box}$ and $C_{HD}$~\cite{Alonso:2013hga}.
Along these lines we expect the mixing of $Q_{q^2H^4D}^{(2)}$ to be closely related to that of $Q_{q^2H^4D}^{(3)}$ because, as shown in Sec.~\ref{sec:cls_ops}, these terms contain orthogonal linear combinations of operators.
A similar relation is expected for $Q_{l^2H^4D}^{(2, 3)}$.

%--------------------------------------------------------------------------------------------------------------
\begin{table}[H]
\begin{center}
\renewcommand{\arraystretch}{1.2}
\begin{tabular}[t]{c c *5{| p{2.35cm}} }
%-----------------
\multirow{5}{*}{$\overline w$} & 8 &
\cellcolor[gray]{0.94}$X_L^4$ & 
\cellcolor[gray]{0.88}$X_L^3 H^2$, ${\color{blue} X_L^2 \psi^2 H}$, ${\color{red} X_L \psi^4}$        & 
\cellcolor[gray]{0.82}${\color{red} X_L^2 H^4}$, ${\color{red} X_L \psi^2 H^3}$, ${\color{red}  \psi^4 H^2}$       &    
\cellcolor[gray]{0.76}${\color{red} \psi^2 H^5}$      & 
\cellcolor[gray]{0.70}${\color{red} H^8}$            \\ \cline{2-7}
%-----------------
& 6 &
           & 
 \cellcolor[gray]{0.94}$X_L^2 H^2 D^2$, $X_L^2 \psi \bar \psi D$, ${\color{blue} X_L \psi^2 H D^2}$, ${\color{red} \psi^4 D^2}$ & 
 \cellcolor[gray]{0.88}${\color{red}  X_L H^4 D^2}$, ${\color{blue} X_L^2 \bar \psi^2 H}$, ${\color{red}  X_L \psi \bar \psi H^2 D}$, ${\color{red}  \psi^2 H^3 D^2}$, ${\color{red}  X_L \psi^2 \bar \psi^2}$, ${\color{red}  \psi^3 \bar \psi H D}$   & 
 \cellcolor[gray]{0.82}${\color{red} H^6 D^2}$, ${\color{red} \psi \bar \psi H^4 D}$, ${\color{red} \psi^2 \bar \psi^2 H^2}$             &  
 \cellcolor[gray]{0.76}$ {\color{red} \bar \psi^2 H^5}$                  \\ \cline{2-7}
%-----------------
& 4 &           
           &                           
           & 
\cellcolor[gray]{0.94}$X_L^2 X_R^2$, ${\color{blue} X_L X_R H^2 D^2}$, ${\color{red} H^4 D^4}$, ${\color{blue} X_L X_R \psi \bar \psi D}$, ${\color{blue} X_R \psi^2 H D^2}$, ${\color{blue} {\color{blue} X_L \bar \psi^2 H D^2}}$, ${\color{red} \psi \bar \psi H^2 D^3}$, ${\color{red} \psi^2 \bar \psi^2 D^2}$ & 
\cellcolor[gray]{0.88}${\color{red}  X_R H^4 D^2}$, ${\color{blue} X_R^2 \psi^2 H}$, ${\color{red}  X_R \psi \bar \psi H^2 D}$, ${\color{red} \bar  \psi^2 H^3 D^2}$, ${\color{red}  X_R \psi^2 \bar \psi^2}$, ${\color{red}  \psi \bar \psi^3 H D}$     & 
\cellcolor[gray]{0.82}${\color{red} X_R^2 H^4}$, ${\color{red} X_R \bar  \psi^2 H^3}$, ${\color{red} \bar \psi^4 H^2}$ \\ \cline{2-7}
%-----------------
& 2 &          
           &                           
           &                         
           & 
\cellcolor[gray]{0.94}$X_R^2 H^2 D^2$, $X_R^2 \psi \bar \psi D$, ${\color{blue} X_R \bar \psi^2 H D^2}$, $ {\color{red} \bar \psi^4 D^2}$ & 
\cellcolor[gray]{0.88}$X_R^3 H^2$, ${\color{blue} X_R^2 \bar \psi^2 H}$, ${\color{red} X_R \bar \psi^4}$ \\ \cline{2-7}
%-----------------     
& 0 &      
           &                           
           &                         
           &                           
           & 
\cellcolor[gray]{0.94}$X_R^4$ \\ \cline{2-7}
%-----------------  
& & 0 & 2 & 4 & 6 & 8 \\
\multicolumn{2}{c}{} & \multicolumn{5}{c}{$w$}
\end{tabular}
\end{center}
\caption{Weight lattice for dimension-8 operators in the SMEFT. 
By the non-renormalization theorem of Ref.~\cite{Cheung:2015aba} a subclass of operators, $j$, can renormalize another subclass, $i$, at one-loop if $w_i \geq w_j$ and $\overline w_i \geq \overline w_j$. 
Visually this prevents transitions down or to the left.
Additionally, if the SMEFT is UV-completed by a weak-coupled and renormalizable theory, then only the operator classes in red can be generated via tree level matching, whereas the black and blue operator classes are only generated at loop level in this scenario.
However, the operator classes in blue can be renormalized by the red tree level operator classes even though the blue operator classes are not generated by these types of UV theories at tree level~\cite{Craig:2019wmo}.}
\label{tab:rge}
\end{table}








